\section{Egyptian}

Egypt has no single creation story. Three priestly centers, Heliopolis, Hermopolis, and Memphis, each developed a full account. They coexisted. All three begin in the same place, the primordial waters, and all three set up an order that must then be held against disorder.

\subsection{The Creation Model}

In the beginning there is only an endless dark ocean. It is not nothing. It is undifferentiated potential, formless and motionless, holding everything that could be.

From these waters the first mound of dry land rises. On the mound a self-made one appears. It is complete in itself and needs no partner. By its own substance it brings forth the first pair, a principle of dry expansive space and a principle of moist receptivity. From that pair come the next pair, the ground beneath and the boundary above.

At first the ground and the boundary are locked together. The principle of space lifts the boundary off the ground and holds them apart, opening the room in which all things can exist. From the ground and the boundary come the ruling powers, the generation that governs the world, including a just king, a worker of magic, and a force of disruption that opposes the order.

What the first one establishes is order itself, the standard of truth and balance. Against it presses disorder, which never fully departs because the dark waters never fully go away. The tradition as a whole is the continual work of holding the order in place against the disorder. The clearest daily image of this is the journey of light across the sky and through the dark below, attacked each night by a great serpent of chaos and reborn each dawn.

\subsection{The Starting State}

\begin{itemize}
 \item \textbf{Nun} is the primordial waters. The dark, infinite, formless, motionless ocean that existed before anything. Not nothing, but undifferentiated potential. Every account begins here.
 \item The \textbf{Benben} is the first mound of dry land that rises from Nun. It is the first place and the first point of order. The creator appears upon it. The pyramid and the obelisk tip are later stone images of the Benben.
 \item Onto the mound the creator brings the first sunrise, the first act of differentiation, light from darkness.
\end{itemize}

\subsection{The Three Creation Theologies}

\textbf{Heliopolis, the Ennead.} The creator is \textbf{Atum}, who self-creates from Nun on the first mound and then produces a family of nine gods by generation.

\begin{longtable}{L{0.16\linewidth}
L{0.4\linewidth} L{0.32\linewidth}}
\toprule
\textbf{Generation} & \textbf{Gods} & \textbf{Domain} \\
\midrule
1 & \textbf{Atum} & The self-created creator \\
2 & \textbf{Shu}, \textbf{Tefnut} & Air, moisture \\
3 & \textbf{Geb}, \textbf{Nut} & Earth, sky \\
4 & \textbf{Osiris}, \textbf{Isis}, \textbf{Set}, \textbf{Nephthys} & Kingship, magic, chaos, mourning \\
\bottomrule
\end{longtable}

Atum produces the first pair on his own, by his own substance. Shu, the air, then lifts Nut, the sky, off Geb, the earth, opening the space for existence. The full \textbf{Ennead} is these nine.

\textbf{Hermopolis, the Ogdoad.} Before any creator there were eight primordial gods, four pairs embodying the qualities of the pre-creation chaos. Each pair is a male and a female form of one quality.

\begin{longtable}{L{0.08\linewidth}
L{0.42\linewidth} L{0.38\linewidth}}
\toprule
\textbf{\#} & \textbf{Pair} & \textbf{Quality} \\
\midrule
1 & \textbf{Nun} and \textbf{Naunet} & The primordial waters \\
2 & \textbf{Heh} and \textbf{Hauhet} & Infinity, endlessness \\
3 & \textbf{Kek} and \textbf{Kauket} & Darkness \\
4 & \textbf{Amun} and \textbf{Amaunet} & Hiddenness, the unseen \\
\bottomrule
\end{longtable}

From the interaction of these eight, the first mound rose and the sun was born, in some versions from a cosmic egg, in others from a lotus on the waters.

\textbf{Memphis, the theology of Ptah.} The most abstract account, preserved on the Shabaka Stone. Here \textbf{Ptah} is the supreme creator, prior even to Atum, and he creates not by physical act but by mind and speech.

\begin{enumerate}
 \item Ptah conceives all things in his \textbf{heart}, the seat of thought and will.
 \item Ptah brings them into being by his \textbf{tongue}, by uttering their names. The word spoken makes the thing real.
 \item Atum and the Ennead are themselves the product of Ptah's heart and tongue. Thought and the creative word stand above the older generative myth.
\end{enumerate}

This creation by the word is sometimes compared to the divine word of later traditions.

\subsection{Maat Versus Isfet}

Creation does not end the threat of chaos. It sets up an order that must be maintained.

\begin{itemize}
 \item \textbf{Maat} is the principle of truth, justice, balance, and cosmic order, personified as a goddess with a feather. It is what the creator establishes and what the king, the gods, and every person must uphold.
 \item \textbf{Isfet} is chaos, disorder, falsehood, and injustice. The ever-present counter-force. The waters of Nun never fully go away, and isfet always presses to return.
\end{itemize}

Egyptian religion as a whole can be read as the continual work of keeping Maat in place against isfet.

\subsection{The Solar Cycle}

The clearest daily image of Maat against isfet is the journey of the sun.

\begin{enumerate}
 \item \textbf{Ra} sails across the sky by day in his solar barque, bringing light and life.
 \item At dusk he descends into the underworld (the \textbf{Duat}) and travels through its twelve hours of night.
 \item Each night the great serpent \textbf{Apep}, the embodiment of chaos, attacks the barque to swallow the sun.
 \item Ra and his crew defeat Apep each night, and Ra is reborn at dawn.
\end{enumerate}

The victory is never final. The battle is fought again every night. This is the cosmic order sustained by repetition.

\subsection{The Parts of the Soul}

The Egyptian person was made of several separable parts, each a real component that had to survive and be cared for. The body had to be preserved because some of these parts depended on it.

\begin{longtable}{L{0.16\linewidth}
L{0.62\linewidth}}
\toprule
\textbf{Part} & \textbf{Meaning} \\
\midrule
\textbf{Ka} & The vital essence, the life-force given at birth. It survives death and needs food offerings. \\
\textbf{Ba} & The personality, shown as a human-headed bird. It can leave the tomb and return to the body at night. \\
\textbf{Akh} & The transfigured, blessed spirit a person becomes after a successful judgment, when Ka and Ba unite. \\
\textbf{Ren} & The name. As long as it is spoken and remembered, the person endures. \\
\textbf{Sheut} & The shadow, always present, tied to protection and to the person's substance. \\
\textbf{Ib} & The heart, the seat of thought, will, and morality. It is weighed at judgment. \\
\bottomrule
\end{longtable}

\subsection{The Duat and the Weighing of the Heart}

The \textbf{Duat} is the underworld through which the sun travels at night and through which the dead must pass. It is full of gates, guardians, and dangers, each needing the right spell or name. The funerary texts are the maps and passwords for the journey.

The central judgment scene is the \textbf{Weighing of the Heart}.

\begin{enumerate}
 \item The deceased is led by \textbf{Anubis} into the \textbf{Hall of Two Truths}.
 \item The deceased recites the negative confession, declaring innocence before 42 assessor gods.
 \item The \textbf{heart} is placed on one pan of a great scale. On the other sits the \textbf{feather of Maat}.
 \item \textbf{Anubis} adjusts the scale and \textbf{Thoth} records the result.
 \item If the heart is as light as the feather, the person is declared true of voice and passes on to become an \textbf{akh} in the Field of Reeds.
 \item If the heart is heavy with wrongdoing, it is thrown to \textbf{Ammit}, the Devourer. This is the second death, true annihilation, the worst Egyptian fate.
\end{enumerate}

\begin{longtable}{L{0.22\linewidth}
L{0.56\linewidth}}
\toprule
\textbf{Figure} & \textbf{Role at Judgment} \\
\midrule
\textbf{Anubis} & Jackal god, guides the dead, oversees the scale \\
\textbf{Thoth} & Ibis-headed scribe, records the verdict \\
\textbf{Maat} & Her feather is the standard of weighing \\
\textbf{Ammit} & Devours the hearts of the guilty \\
\textbf{Osiris} & Lord of the dead, presides over the hall \\
\bottomrule
\end{longtable}

\subsection{Comparison}

The Egyptian model fits the base shape but stresses one part of it more than the others. The dark waters give a self-made first one, the first one differentiates into paired powers, and a standard of order is set. Where Norse renewal is built into a single great ending, here the threat of disorder is daily and the order is held by repetition rather than by one final battle. The pattern still runs through it. \textbf{Undifferentiated potential differentiates into a structured world that must be actively held against the disorder it came from.}
